\section{Priority boundaries and interpretation}
The real cone identity in Lemma~\ref{lem:balance} is not new: quadratic rank-one generation and semidefinite representations for one quadratic constraint have a substantial literature, including Sturm--Zhang \cite{SZ}. Nishijima--Louren\c{c}o \cite{NL} classify real faces of copositive and completely positive cones over a second-order cone. Their real facial classification is an important comparator, not something superseded merely by using arithmetic terminology.

Rational interior reconstruction has classical antecedents, including rational completely positive factorization \cite{DSV}; the present reserve proof is included to make the precise density hypothesis transparent. Rational orthogonal density, signed Cayley parametrizations, and quantitative rational approximation also predate this work \cite{Schmutz,LO,TWZ}. The exact width theorem is a consequence of these facts and positive-scaling openness, not a new parametrization of the orthogonal group. The $r+3$ unrestricted bound is derived from Meyer's theorem and is credited as classical quadratic-form arithmetic.

The Hasse--Minkowski principle, Meyer threshold, and small-zero bounds are classical \cite{CasselsBook,CasselsSmall,Dietmann}. The candidate new content is their exact integration into constrained rank-one generation and arbitrary affine Gram completion, including the maximal-rational-support step, the single-isotropy-test carrier, and the sharp rational spherical-moment obstruction. The product-cone circuit reduction supplies an explicit scope boundary. Burer--Dong \cite{BD} already represent nonconvex QCQPs through generalized completely positive cones over products of second-order cones and over PSD cones. Their bounded real convexification theorem and proof have been inspected; the present formulation isolates exact rational feasibility, only two homogeneous observations, unrestricted widths, and a binary-height lower bound without a boundedness promise. We do not claim that conic universality itself is new. We have not established that no earlier publication contains an equivalent formulation.

The earlier qualitative generalized-perfect-matrix theory \cite{OS} is also not a new result here. No blanket rational boundary factorization theorem is asserted for completely positive matrices or arbitrary rational cones. In particular, the Lorentz counterexamples here concern a different column cone and do not depend on reported 2026 orthant counterexamples.

\paragraph{Publication status.} The prior registered LAA comparison has not been completed or scored by this continuation. No numerical journal percentile, first-priority certificate, acceptance probability, external human verification, or proof-assistant verification is asserted. The contribution should be assessed primarily by the complete arithmetic classification and its affine and moment consequences, not by the number of examples or the name of a target venue.

\begin{thebibliography}{99}
\bibitem{SZ} J. F. Sturm and S. Zhang, On cones of nonnegative quadratic functions, \emph{Mathematics of Operations Research} 28 (2003), 246--267. DOI: 10.1287/moor.28.2.246.14485.
\bibitem{BD} S. Burer and H. Dong, Representing quadratically constrained quadratic programs as generalized copositive programs, \emph{Operations Research Letters} 40 (2012), 203--206. DOI: 10.1016/j.orl.2012.02.001.
\bibitem{NL} M. Nishijima and B. F. Louren\c{c}o, Facial structure of copositive and completely positive cones over a second-order cone, \emph{Linear Algebra and its Applications} 743 (2026), 211--249. DOI: 10.1016/j.laa.2026.04.019. Preprint arXiv:2502.04006v2.
\bibitem{SEDZ} M. Safey El Din and L. Zhi, Computing rational points in convex semi-algebraic sets and sum of squares decompositions, \emph{SIAM Journal on Optimization} 20 (2010), 2876--2889. DOI: 10.1137/090772459. In particular Theorem 1.1 and the real-algebraic sampling preliminaries.
\bibitem{CasselsBook} J. W. S. Cassels, \emph{Rational Quadratic Forms}, London Mathematical Society Monographs 13, Academic Press, 1978.
\bibitem{CasselsSmall} J. W. S. Cassels, Bounds for the least solutions of homogeneous quadratic equations, \emph{Proceedings of the Cambridge Philosophical Society} 51 (1955), 262--264.
\bibitem{Fukshansky} L. Fukshansky, Heights and quadratic forms: Cassels' theorem and its generalizations, \emph{Contemporary Mathematics} 587 (2013), 77--94. DOI: 10.1090/conm/587/11673.
\bibitem{Dietmann} R. Dietmann, Small zeros of quadratic forms avoiding a finite number of prescribed hyperplanes, \emph{Canadian Mathematical Bulletin} 52 (2009), 63--65. DOI: 10.4153/CMB-2009-007-7.
\bibitem{Schmutz} E. Schmutz, Rational points on the unit sphere, \emph{Central European Journal of Mathematics} 6 (2008), 482--487. DOI: 10.2478/s11533-008-0038-4.
\bibitem{LO} H. Liebeck and A. Osborn, The generation of all rational orthogonal matrices, \emph{American Mathematical Monthly} 98 (1991), 131--133.
\bibitem{TWZ} Q. Tang, W. Wang and H. Zhang, Generating all regular rational orthogonal matrices, \emph{Linear Algebra and its Applications} (2025). DOI: 10.1016/j.laa.2025.06.008. Preprint arXiv:2410.11246, whose Theorem 3 recalls Liebeck--Osborn.
\bibitem{DSV} M. Dutour Sikiri\'c, A. Sch\"urmann and F. Vallentin, Rational factorizations of completely positive matrices, \emph{Linear Algebra and its Applications} 523 (2017), 46--51. DOI: 10.1016/j.laa.2017.02.017.
\bibitem{OS} A. Oertel and A. Sch\"urmann, Generalized Perfect Matrices, arXiv:2602.05841v2 (2026).
\end{thebibliography}
